\bookchapter{פרק 13: המלצות על אפליקציות כלים}

\chapterintro{בפרק זה בחרנו כלי פרודוקטיביות, בריאות ונסיעות, והשלמנו יכולות מערכת שימושיות ויכולות בין-מכשירים: ניקוי וחסימת שיחות זבל כמו Qingzhuo, מתגי חסימה מובנים בטלפון, התראות מזג אוויר ורעידות אדמה; שליטה מרחוק במחשב דרך Sunlogin/ToDesk, עבודה עם AI כמו Trae, שולחן עבודה מלאכותי, העברת קבצים מהירה ברשת מקומית כמו LocalSend, לוח-גזירים בין-מכשירים, קודי אימות דינמיים כמו Aegis/Microsoft Authenticator, לוח שנה ומשימות של TickTick וסריקת נייר.}

\booksection{כלי פרודוקטיביות: ניהול הזמן והקשב שלך}

כלי פרודוקטיביות הם “השלד” של החיים הדיגיטליים — הם עוזרים לארגן רעיונות, משימות וחומרים מפוזרים בצורה מסודרת. להלן המלצות לפי תרחיש.

\booksubsection{רשימת משימות (ניהול משימות)}

\begin{itemize}
  \item TickTick — למי מתאים: למי שצריך לרכז במקום אחד את משימות העבודה, החיים והלימודים יחד עם הלוח השנה; אחת ההמלצות המרכזיות של הפרק. נקודות החוזקה העיקריות: סנכרון בין נייד, מחשב, טאבלט ו-Web; רשימות + לוח שנה + טיימר פומודורו בחבילה אחת; הזנה בשפה טבעית — "פגישה מחר ב-3 אחר הצהריים", "דוח שבועי כל יום שני" — עם זיהוי אוטומטי של זמנים וכללי חזרה; תמיכה בתתי-משימות, תגיות, עדיפויות, קבצים מצורפים ורשימות שיתופיות. מתאים לניהול מאוחד של "תכנון יומן + סימון משימות + סקירה יומית".
  \item איך להשתמש ב-TickTick בצורה חלקה יותר: תיבת איסוף: תחילה זרקו הכול ל"היום/תיבת איסוף", ובערב סיווגו לרשימות עבודה, משפחה ולימודים.
  \item תצוגת לוח שנה: גררו פגישות ותאריכי יעד ללוח השנה כדי לראות התנגשויות — אינטואיטיבי יותר מרשימת משימות בלבד.
  \item משימות חוזרות: תרופות, תשלומים ואימונים הוגדרו כמשימות עם כללים — לא צריך למלא מחדש בכל פעם.
  \item טיימר פומודורו: התחילו מיקוד ישירות על משימה, וספרו כמה "באמת סיימתם" השבוע.
  \item שילוב עם פונקציות הטלפון: העוזרת הקולית יכולה להוסיף תזכורות ב"הזכירי לי…"; לתאריכי יעד חשובים הפעילו גם לוח שנה במערכת כפול-ביטוח.
\end{itemize}

Todoist קהל יעד: משתמשים שמחפשים מינימליזם ועבודה בין פלטפורמות. נקודות חוזק: ממשק נקי ומאופק, היררכיית פרויקטים ברורה (משימה תת-משימה תת-תת-משימה), נקודות karma לעידוד הרגלים.

Microsoft To Do קהל יעד: עובדי משרד ומשתמשי Outlook. נקודות חוזק: חינם לחלוטין, בלי פרסומות, "היום שלי" מתרוקן כל יום ומפחית עומס נפשי, אינטגרציה עמוקה עם מיקרוסופט.

פתקים/תזכורות במערכת קהל יעד: משתמשים קלים שלא רוצים להתקין אפליקציות צד שלישי. נקודותחוזק: חווייה מובנית, אפס עיכובים, הוספה קולית דרך Siri/XiaoAi, שילוב עם ווידג'טים של המערכת.

\booksubsection{טיימר פומודורו (אימון קשב)}

\begin{itemize}
  \item Forest קהל יעד: מי שלא שולט ביד וגולש בקלות בטלפון. נקודות חוזק: שותלים עץ בזמן מיקוד, אם יוצאים מוקדם העץ מת — בונים משוב חיובי דרך "חיים וירטואליים", אפשר לשתול עם חברים.
  \item Tide — למי מתאים: למי שזקוק לרעש לבן כדי להתרכז או להירדם. נקודות החוזקה העיקריות: צלילי טבע (גשם, גלי ים, יער) + פומודורו + מדיטציה בחבילה אחת, עם ממשק מינימליסטי ומרגיע במיוחד.
  \item עקרון: 25 דקות מיקוד קלאסיות + 5 דקות מנוחה, אפשר להתאים ל-50/10 או 90/20 לפי המצב.
\end{itemize}

\booksubsection{לוח שנה (ויזואליזציה של זמן)}

\begin{itemize}
  \item לוח שנה של המערכת קהל יעד: כל המשתמשים. נקודותחוזק: התראות מובניות יציבות ביותר, צריכת סוללה נמוכה מאוד, הוספה דרך Siri/XiaoAi.
  \item Google Calendar קהל יעד: משתמשים עם כמה מכשירים, ומי שצריך לשתף לוחות זמנים. נקודותחוזק: "מטרות" מתכננות אוטומטית זמן לאימון/לימוד, זיהוי חכם של אירועים מהודעות Gmail.
\end{itemize}

\booksubsection{פתקים (צבירת ידע)}

\begin{itemize}
  \item Notion — למי מתאים: למי שצריך לנהל במקום אחד פתקים, משימות, מסדי נתונים וויקי; סטודנטים, פרילנסרים וצוותים. נקודות החוזקה העיקריות: בסיס ידע כל-יכול עם עריכה מודולרית בבלוקים (Blocks), אפשרות לבנות דף בית אישי, רשימות קריאה ולוחות פרויקטים, ואקוסיסטם תבניות עשיר.
  \item Obsidian קהל יעד: מנהלי ידע מעמיקים, חוקרים וכותבים. נקודותחוזק: אחסון Markdown מקומי, קישורים דו-כיווניים + גרף קשרים — הפתקים שלך הופכים ל"מוח שני", הנתונים שלך בלבד.
  \item flomo קהל יעד: מי שרוצה לתעד רעיונות מזדמנים בלי לכתוב מאמרים ארוכים. נקודותחוזק: פתקי כרטיסים, כמו לפרסם Moments, תמיכה בתגיות \#, חזרה לפי ציר זמן.
  \item Apple Notes קהל יעד: משתמשי iPhone/Mac עם צורכי תיעוד קלים. נקודותחוזק: פתיחה מהירה, סריקת מסמכים/ציור ידני/רשימות/סנכרון iCloud, חינם ואינטגרציה עמוקה עם המערכת.
  \item Evernote קהל יעד: משתמשים מתקדמים שצריכים הרבה גזירי אינטרנט וארגון חומרים. נקודותחוזק: תוסף גזירת אינטרנט חזק, OCR לזיהוי טקסט בתמונות, ותיק ויציב, אבל הגרסה המתקדמת בתשלום.
\end{itemize}

\booksubsection{כלי PDF}

\begin{itemize}
  \item PDF Expert קהל יעד: משתמשי iPad/iPhone, אנשי עסקים שצריכים להעיר על חוזים, סטודנטים. נקודותחוזק: עורך ה-PDF החזק ביותר ב-iOS, סימון חלק, חתימות, מיזוג/פיצול, מילוי טפסים.
  \item WPS קהל יעד: עובדי משרד שעובדים עם Word/Excel/PPT/PDF. נקודותחוזק: חבילת משרד כול-באחרת, תאימות טובה לכונני ענן בסין ולקבצי WeChat, הגרסה החינמית מספיקה ליום-יום.
  \item Adobe Acrobat קהל יעד: מי שדורש עמידה בתקני PDF (משפט, עיצוב, הוצאה לאור). נקודותחוזק: תקן רשמי, המרה ל-PDF/A מדויקת, סנכרון בין מכשירים וחתימה אלקטרונית מקצועיים.
\end{itemize}

\booksubsection{סריקה (דיגיטציה של נייר)}

חוזי נייר, חשבוניות, שיעורי בית, פתקי לוח, צילומי תעודות… מצלמת הטלפון כבר מחליפה את רוב הסורקים. תהליך ליבה: לצלם ברור חיתוך ותיקון אוטומטי ייצוא PDF/תמונה OCR לפי בחירה.

\begin{itemize}
  \item CamScanner קהל יעד: אנשי עסקים וסטודנטים שסורקים הרבה חוזים, קבלות ותעודות. נקודותחוזק: זיהוי קצה אוטומטי, תיקון פרספקטיבה, OCR, אפשר ליצור PDF מרובה עמודים ולסנכרן לענן.
  \item Microsoft Lens קהל יעד: סטודנטים, מי שצריך לסרוק לוח/עמוד ספר ל-Word. נקודותחוזק: חינם בלי פרסומות, ייצוא בלחיצה ל-OneNote/Word, דיוק גבוה בזיהוי טקסט.
  \item סריקה ב-Apple Notes קהל יעד: משתמשי iPhone, צורכי סריקה קלים. נקודותחוזק: בלי להתקין אפליקציה, לוחצים על סמל המצלמה בפתקים וסורקים, הסרת צללים אוטומטית, חתימה וסימון, חינם ונוח.
  \item Baimiao / גרסת CamScanner באינטרנט / "סריקת מסמכים" במצלמת המותג: מצלמות המערכת של Huawei, Xiaomi, OPPO, vivo כוללות לרוב מצב "מסמך/סריקה"; ב-Android אפשר גם Google Camera או סריקת המערכת. מתאים לצילום דחוף של עמוד חוזה.
  \item טיפים לצילום: אור אחיד, הטלפון מקביל לנייר, כוונו לארבע הפינות; קבצים מרובי עמודים צלמו לפי הסדר ואז מזגו ל-PDF; בתעודות רגישות שימו לב ל-watermark ולמיקום השמירה, ומחקו לחלוטין מ"נמחק לאחרונה" אחרי השימוש.
\end{itemize}

\booksection{בריאות וכושר}

אפליקציות בריאות הן “העוזר האישי” הכי מוערך בטלפון — הן עוקבות ברצף אחרי שינה, דופק ופעילות, ומשקפות מצב ארוך-טווח טוב יותר מביקור חד-פעמי בבית חולים.

\booksubsection{כושר וספורט}

\begin{itemize}
  \item Keep קהל יעד: מתאמנים בבית, מתחילים מאפס, מי שרוצה לרדת במשקל ולחטב. נקודותחוזק: המון קורסים חינם, מ-HIIT ועד יוגה, תוכנית אימון אוטומטית מבוססת AI, קהילה חזקה.
  \item Nike Training Club (NTC) קהל יעד: מי שמחפש הדרכת מאמן מקצועי ומערכת אימון מסודרת. נקודותחוזק: רשמי של Nike, כל הקורסים חינם, הדגמות על ידי מאמנים מקצועיים, סיווג ברור לפי קבוצות שרירים/ציוד.
  \item Apple Fitness קהל יעד: משתמשי Apple Watch, מי שרוצה מוטיבציה דרך טבעות סגירה. נקודותחוזק: שלוש טבעות יומיות (תנועה/אימון/עמידה), מטרות ויזואליות, אינטגרציה חלקה עם נתוני דופק מהשעון.
\end{itemize}

\booksubsection{מעקב שינה}

\begin{itemize}
  \item AutoSleep קהל יעד: משתמשי Apple Watch, מי שרוצה לדעת את יחס השינה העמוקה/קלה בכל לילה. נקודותחוזק: מעקב אוטומטי בשינה עם שעון, בלי הפעלה/עצירה ידנית, ציון איכות שינה והמלצות לשינה.
  \item Sleep Cycle קהל יעד: מי שלא ישן עם שעון ורוצה לנטר מהטלפון. נקודותחוזק: ניתוח נחירות/התהפכות דרך מיקרופון, השכמה חכמה בשלב שינה קלה, קימה בלי סבל.
  \item Mi Fitness קהל יעד: משתמשי צמיד/שעון Xiaomi. נקודותחוזק: חינם ומשלב עמוק עם צמיד Xiaomi, מעקב משולב של דופק/חמצן בדם/שינה, יחס מחיר-ביצועים מעולה.
\end{itemize}

\booksubsection{מדיטציה והרפיה}

\begin{itemize}
  \item Calm קהל יעד: עובדים תחת לחץ, מתקשים להירדם. נקודותחוזק: מדיטציה מודרכת "שלווה יומית" + סיפורי שינה (בקריאה של מפורסמים), ממשק יפה, תוכן אנגלי עשיר.
  \item Tide — למי מתאים: למשתמשים בסינית שזקוקים לרעש לבן לריכוז או לשינה. נקודות החוזקה העיקריות: מוצר סיני משובח, צלילי טבע + קורסי מדיטציה + פומודורו, עם הדרכה בסינית שמרגישה קרובה ואישית יותר.
  \item Xiao Sleep — למי מתאים: למי שצריך צלילים לשינה טובה ורוצה לתעד דיבור בשינה או נחירות. נקודות החוזקה העיקריות: מאגר עצום של צלילי רעש לבן שאפשר לערבב ולשלב (גשם + מאוורר + מדורה), כיבוי מתוזמן לפני השינה, והקלטת דיבור בשינה — משעשע ושימושי כאחד.
\end{itemize}

\booksubsection{בריאות האישה}

\begin{itemize}
  \item Flo קהל יעד: מי שצריכה לתעד מחזור, תכנון הריון או בריאות האישה. נקודותחוזק: חיזוי AI מדויק למחזור הבא/ביוץ, קהילה עם שאלות אנונימיות, עיצוב ממשק מוקפד.
  \item Meiyou — למי מתאים: למשתמשות בסין ולמי שמעוניין בקהילה ובתוכן מדעי פופולרי בסינית. נקודות החוזקה העיקריות: תוכן מקומי עשיר, עם ידע מותאם לכל שלב — מחזור, ניסיון להרות, היריון או גידול ילדים.
\end{itemize}

\booksubsection{תזכורות לתרופות}

\begin{itemize}
  \item Yongyao Zhushou (עוזר תרופות) — למי מתאים: לחולים כרוניים שנוטלים תרופות לטווח ארוך, ולבני משפחה שמסדרים את התרופות של הורים מבוגרים. נקודות החוזקה העיקריות: זיהוי תרופות בסריקת ברקוד ובנייה אוטומטית של רשימת תרופות, תזכורת קופצת בזמן, ותיעוד הנטילה שאפשר לשתף עם בני המשפחה.
\end{itemize}

\booksubsection{מרכז נתונים}

\begin{itemize}
  \item Apple Health / Google Fit קהל יעד: כל משתמשי הסמארטפון. נקודותחוזק: מרכז את כל הנתונים מהטלפון, הצמיד ואפליקציות בריאות צד שלישי (צעדים, דופק, שינה, מחזור ועוד) — "תיק בריאות" מאוחד.
\end{itemize}

\booksection{נסיעות ותנועה}

בחוץ, כמה אפליקציות טובות חוסכות המון טעויות. להלן המלצות לפי שלבי הנסיעה.

\booksubsection{טיסות ורכבות מהירות}

\begin{itemize}
  \item Umetrip — למי מתאים: לאנשי עסקים שטסים לעיתים קרובות. נקודות החוזקה העיקריות: נתוני מנהל התעופה האזרחית הרשמיים — המקור המוסמך ביותר לסטטוס טיסות, שינויי שערים והתראות עיכוב; אפשר לבדוק גם דגם מטוס ומפת מושבים.
  \item 12306 קהל יעד: כל מי שנוסע ברכבת בסין. נקודותחוזק: ערוץ הרכישה הרשמי היחיד האמין, תמיכה ברכישת המתנה (קנייה אוטומטית כשיש ביטולים), נוח לכרטיסי סטודנט/קבוצות.
  \item VariFlight — למי מתאים: למי שחשובה לו העמידה בזמנים של טיסות ולמי שצריך להסיע לשדה התעופה או לאסוף ממנו. נקודות החוזקה העיקריות: חיזוי עיכובי טיסות מדויק, ושיתוף המסלול עם בני משפחה שיראו בזמן אמת את מיקומך ואת שעת ההגעה.
\end{itemize}

\booksubsection{מפות וניווט}

\begin{itemize}
  \item Amap — למי מתאים: למשתמשים בערים בסין, לנהגים ולנוסעי מוניות. נקודות החוזקה העיקריות: תנועה בזמן אמת מדויקת ויכולת חזקה לעקוף פקקים, הזמנת מונית ונסיעה משותפת במקום אחד, ומגוון עשיר של מסלולי החלפה בתחבורה ציבורית.
  \item Baidu Maps קהל יעד: מי שזקוק למידע על מקומות וניווט פנימי. נקודותחוזק: מפות פנורמה/ניווט AR, מפות פנים מפורטות לאטרקציות/קניונים, חבילות קול עשירות.
  \item Google Maps קהל יעד: מטיילים בחו"ל. נקודותחוזק: הכיסוי העולמי הרחב ביותר, ביקורות מקומיות אמיתיות, Street View ומפות לא מקוונות — הצלה בארץ זרה.
\end{itemize}

\booksubsection{תרגום}

\begin{itemize}
  \item Google Translate קהל יעד: מטיילים בחו"ל, לומדי שפות. נקודותחוזק: "שיחה" עם תרגום דו-כיווני בזמן אמת, "מצלמה" לתרגום מיידי של מילים — כלי מושלם לטיול, תמיכה בחבילות שפה לא מקוונות.
\end{itemize}

\booksubsection{המרת מטבע}

\begin{itemize}
  \item שערי חליפין מינימליים (Jijian Huilü) — למי מתאים: לנוסעים לחו״ל, לקונים מחו״ל ולקניות חוצות גבולות. נקודות החוזקה העיקריות: ממשק נקי במיוחד, תמיכה ביותר מ-160 מטבעות, עדכון שערים בזמן אמת, השוואה של כמה מטבעות במקביל והמרה מהירה בלחיצה ארוכה על המספר.
\end{itemize}

\booksubsection{מזג אוויר}

\begin{itemize}
  \item מזג אוויר של המערכת / רשת מזג האוויר הסינית קהל יעד: כולם, בחירה ראשונה לתחזית והתראות רשמיות. נקודותחוזק: אפליקציות היצרן מתחברות לערוצים רשמיים כמו "מזג האוויר של סין" — התראות על גשם כבד, טייפון, חום גבוה מגיעות כבאנרים/התראות, סמכותיות גבוהה.
  \item Caiyun Weather — למי מתאים: למי שצריך תחזית גשם מדויקת עד רמת הדקה. נקודות החוזקה העיקריות: דיוק מוביל בתעשייה בשאלה "האם ירד גשם בשעתיים הקרובות", מפה דינמית של חזית הגשם שקל להבין במבט, והתראה יזומה לפני שהגשם מתחיל.
  \item Carrot Weather קהל יעד: חובבי מזג אוויר שמחפשים כיף וווידג'טים עשירים. נקודותחוזק: AI עם קול ציני שמשדר מזג אוויר, המון ווידג'טים מותאמים, מפת מכ"ם מקצועית, הגרסה בתשלום עשירה יותר.
  \item עיקרי התראות: אפליקציות מזג אוויר חייבות לקבל אישור התראות כדי לקבל התראות; ההתראות בארבע רמות — כחול, צהוב, כתום, אדום — אדום הוא הסיכון הגבוה ביותר למזג אוויר אסון, יש להתאים נסיעות לפי ההנחיות. להפעלה נוספת ראו סעיף 7 בפרק "התראות מזג אוויר / אסונות טבע / רעידות אדמה".
\end{itemize}

\booksubsection{מדריכי טיול}

\begin{itemize}
  \item Mafengwo — למי מתאים: לאוהבי תכנון מסלולים מפורט ולטיולי רכב עצמאי או טיולי עומק. נקודות החוזקה העיקריות: יומני מסע עשירים שנכתבו על ידי משתמשים, מדריכים מפורטים לפי יעד — פעילויות, מסלולים ואיך להימנע ממלכודות תיירות — ואפשרות להוריד מדריכים לשימוש לא-מקוון.
  \item Xiaohongshu — למי מתאים: לטיילים צעירים שרוצים חוויות אותנטיות ומחפשים פינות חמד או מקומות אוכל ייחודיים. נקודות החוזקה העיקריות: תוכן ויזואלי של תמונות וסרטונים קצרים, עדכונים מהירים ועדכניים של תכנים כמו "מסלול 3 ימים בעיר X", וחוויית חיפוש מעולה.
\end{itemize}

\booksection{מבוא לכלי אוטומציה}

אוטומציה הוא מנוף של “הגדר פעם אחת, חסוך לטווח ארוך” — אפשר לצמצם עשרות פעולות חוזרות ביום לאפס. בפרק זה רק מבוא למושג; למשחקים מפורטים ומדריכים ראו פרק 23: חשיבה תכנותית ואוטומציה.

\begin{itemize}
  \item "קיצורי דרך" ב-iOS: כלי אוטומציה רשמי של Apple, בלי תכנות — אפשר להשלים בלחיצה אחת רצף פעולות כמו "להגיע הביתה: הפעלת WiFi + הורדת בהירות + השמעת תחזית מזג אוויר".
  \item אוטומציה ב-Android: קישור סצנות ברמת המערכת (כמו "משימות אוטומטיות" של Xiaomi, "חוכמה לפי מצב" של Huawei) — הכי קל להתחיל; Tasker חזק אבל עקומת למידה תלולה, מתאים לגיקים.
  \item חשיבה ליבה: קודם לזהות "פעולות שחוזרות כל יום", ואז לחשוב "האם אפשר להחליף בטריגר אחד + כמה פעולות". דוגמאות מפורטות ומשחקים מתקדמים ראו בפרק 23.
\end{itemize}

\booksection{דפדפנים ומנועי חיפוש: שער המידע שלך}

הדפדפן הוא החלון שלך לאינטרנט, ומנוע החיפוש הוא המצפן שלך בים המידע. שימוש נכון בהם חוסך המון עיקופים.

\booksubsection{איך עובד דפדפן:}

אחרי שאתה מזין כתובת בדפדפן: (1)הדפדפן שואל שרת DNS על כתובת ה-IP של הדומיין; (2)שולח בקשת HTTP/HTTPS לשרת בכתובת הזו; (3)השרת מחזיר את קוד העמוד (HTML, CSS, JavaScript); (4)הדפדפן מנתח את הקוד ומציג את הדף שאתה רואה.

\booksubsection{המלצות לדפדפנים מובילים:}

\begin{itemize}
  \item Microsoft Edge: מבוסס על ליבת Chromium, מהיר ותואם היטב. כולל עוזר Copilot AI שיכול לסכם עמודים, לתרגם ולשוחח. תומך בסנכרון בין מכשירים — מומלץ מאוד כדפדפן ראשי.
  \item Safari: הדפדפן המוגדר כברירת מחדל במכשירי Apple, אינטגרציה עמוקה עם המערכת, חיסכון באנרגיה ויעיל. פרטיות טובה, אבל יש פחות תוספים.
  \item Chrome: שיעור השימוש הגבוה ביותר בעולם, תוספים עשירים. אבל צורך הרבה זיכרון, ובמובייל צורך יותר סוללה.
\end{itemize}

\booksubsection{בחירת מנוע חיפוש:}

\begin{itemize}
  \item Bing: מנוע החיפוש של Microsoft, אינטגרציה עמוקה עם Edge. תקציר AI נותן תשובה ישירות ולא רק קישורים — יעילות חיפוש גבוהה. חוויית חיפוש תמונות גם טובה.
  \item Yandex: מנוע החיפוש הגדול ברוסיה, יכולת חזקה בחיפוש תמונות וחיפוש תמונה הפוך. שימושי מאוד למציאת תמונות דומות וזיהוי עצמים.
  \item Baidu: הכלי המרכזי לחיפוש מידע בסין, אבל יש הרבה פרסומות — צריך יכולת הבחנה.
  \item Google: מנוע החיפוש המדויק ביותר בעולם, אבל בסין צריך רשת מיוחדת כדי להשתמש בו.
\end{itemize}

\booksubsection{מדריך לאחזור מידע יעיל:}

\begin{itemize}
  \item נצלו חיפוש AI: תקציר AI של Bing ו-Copilot של Edge עונים ישירות על שאלות ומחלצים מידע מרכזי מתוך המון תוצאות — יעיל הרבה יותר מלחיצה על קישור אחר קישור.
  \item זיהוי מידע כוזב ופרסומות: פריטים שמסומנים "פרסומת" או "ממומן" בתוצאות החיפוש הם קידום בתשלום, לא תוצאות אורגניות. שימו לב לסמכות המקור — העדיפו אתרים רשמיים ומדיה מוכרת.
  \item הורדת אפליקציות נכון: הקלידו "שם האפליקציה + אתר רשמי", לדוגמה "WeChat אתר רשמי", והעדיפו תוצאה עם דומיין רשמי. אל תתקינו מאתרי הורדות צד שלישי — לעתים קרובות הם מצרפים תוכנות לא רצויות.
  \item איך לזהות אתר רשמי: להורדת Windows 11 גשו ל-microsoft.com, להורדת Chrome ל-google.com. לבדוק שהדומיין תואם למותג — זו שיטת הזיהוי הבסיסית ביותר.
\end{itemize}

\booksubsection{ידע על דומיינים:}

הדומיין הוא “מספר הבית” של האתר — הבנה שלו עוזרת לשפוט את אופי האתר ומקורו:

\begin{itemize}
  \item משמעויות סיומות נפוצות: .com — ארגון מסחרי, הנפוץ ביותר
  \item .cn — דומיין לאומי של סין
  \item .org — ארגון ללא מטרות רווח
  \item .net — ארגון שירותי רשת
  \item .edu — מוסד חינוכי
  \item .gov — מוסד ממשלתי
  \item .xyz/.tech/.info — דומיינים כלליים חדשים, סף כניסה נמוך — יש להיזהר
\end{itemize}

דומיינים לאומיים/אזוריים: .jp — יפן

.kr — קוריאה

.uk — בריטניה

.de — גרמניה

.tw — טאיוואן

.hk — הונג קונג

אבטחה: אתרים עם סיומת .gov ו-.edu אמינים יותר בדרך כלל, כי סף הבקשה גבוה. דומיינים זולים כמו .xyz ו-.top קלים לניצול על ידי אתרי פישינג — יש להיזהר בגלישה.

\booksubsection{חסימת פרסומות:}

\begin{itemize}
  \item AdGuard קהל יעד: כל מי שסובל מחלונות קופצים, פרסומות וידאו ופרסומות בתוך אפליקציות. נקודות חוזק: חסימת פרסומות ברמת המערכת (לא רק בדפדפן), סינון באנרים ומעקב באפליקציות, חיסכון בתעבורה והאצת טעינה; הגרסה בתשלום תומכת בכל המכשירים.
\end{itemize}

\booksubsection{דפדפני פרטיות:}

\begin{itemize}
  \item Brave קהל יעד: גולשים ששמים דגש על פרטיות ולא רוצים מעקב פרסומי. נקודות חוזק: מבוסס Chromium, חוסם פרסומות ועוקבים כברירת מחדל, מצב Tor מובנה, מהיר יותר מ-Chrome עם פחות זיכרון.
  \item DuckDuckGo קהל יעד: מי שרוצה "לחפש מילה אחת בלי שכל הרשת תרדוף אחריו עם פרסומות". נקודותחוזק: מנוע חיפוש מובנה באותו שם (שומר היסטוריית חיפוש), ניקוי כל העקבות בלחיצה אחת ("כפתור השריפה"), חווייה אחידה בין פלטפורמות.
\end{itemize}

\booksubsection{השוואת דפדפנים:}

\begin{center}
\begin{tabularx}{\linewidth}{|X|X|X|}
\hline
\textbf{דפדפן} & \textbf{נקודות חוזק} & \textbf{קהל יעד} \tabularnewline
\hline
Microsoft Edge & ליבת Chromium + Copilot AI מובנה, מסכם עמודים/מתרגם/משוחח & עובדי משרד שמחפשים יעילות, משתמשי AI כבדים (מומלץ כראשי) \tabularnewline
Safari & אינטגרציה עמוקה עם Apple, חיסכון באנרגיה, פרטיות חזקה & משתמשי iPhone/Mac, שמים דגש על סוללה וקישוריות מערכת \tabularnewline
Chrome & שיעור השימוש הגבוה ביותר בעולם, מערכת התוספים העשירה ביותר & משתמשים מתקדמים שצריכים הרבה תוספים, אוהבי סנכרון בין מכשירים \tabularnewline
Brave & חוסם פרסומות ועוקבים כברירת מחדל, מהיר יותר & מי ששם דגש על פרטיות ורוצה לחסום פרסומות בלי התקנת תוספים \tabularnewline
DuckDuckGo & מנוע חיפוש מובנה לא שומר היסטוריה, ניקוי עקבות בלחיצה אחת & משתמשים רגישים לפרטיות שלא רוצים מעקב מותאם אישית \tabularnewline
\hline
\end{tabularx}
\end{center}

\begin{tcbpractice}{נסו ליצור תהליך אוטומציה "הירגעות לפני שינה"}
\begin{enumerate}
  \item יצירת טריגר: שעה מסוימת (למשל 22:30).
  \item הוספת פעולות: הורדת עוצמת שמע > הפעלת מצב "לא להפריע" > השמעת רשימת הירגעות > הורדת בהירות.
  \item שמרו ובדקו.
\end{enumerate}

\end{tcbpractice}

\booksection{כלי מערכת שימושיים: ניקוי, חסימת שיחות זבל ואבטחה}

בנוסף לאפליקציות פרודוקטיביות ובריאות, מה שבאמת משפיע על החווייה היומיומית הם שלושה דברים: “פחות הפרעות, מספיק אחסון, נוכלים לא נכנסים”. בסעיף זה נסקור תוכנות שימושיות נפוצות, ואז נרחיב בנפרד על חסימת שיחות הזבל המובנית בטלפון — אנשים רבים קנו טלפון לפני שנים ומעולם לא פתחו את המתגים האלה.

\booksubsection{ניקוי מערכת ואופטימיזציית אחסון}

\begin{itemize}
  \item Qingzhuo — למי מתאים: למשתמשי Android שרוצים לנקות בקלילות מטמון, קבצים שנשארו ותיקיות ריקות, בלי להיגרר לחבילות "מנהל הטלפון" המסחריות. נקודות החוזקה העיקריות: נפח קטן ואפס פרסומות, סריקה עמוקה של שאריות אפליקציות, איתור קבצים גדולים וניתוח ניצול האחסון; מרוסן יותר מאפליקציות "האצה בלחיצה אחת" המסחריות, ומתאים למשתמשים מיומנים לתחזוקה תקופתית. שימו לב: לפני הניקוי ודאו את הנתיבים — אל תסמנו באופן עיוור את הכול בתיקיות של WeChat, אלבום התמונות והמסמכים; לניקוי מטמון ברמת המערכת עדיף להשתמש קודם כל בכלי הניקוי המובנה של היצרן, וכלים של צד שלישי רק כתוספת.
  \item ניקוי אחסון מובנה במערכת קהל יעד: כולם, במיוחד קשישים ומי שמפחד מתקלות. נקודותחוזק: Android [הגדרות] > [אחסון] > [ניקוי והאצה]; iOS [הגדרות] > [כללי] > [אחסון iPhone]. המערכת מכירה את מבנה המטמון של האפליקציות שלה — סיכון המחיקה בשוגג הנמוך ביותר, וצריכה להיות הבחירה הראשונה.
  \item ניהול אחסון מובנה ב-WeChat/QQ קהל יעד: משתמשים עם הרבה צ'אטים. נקודותחוזק: WeChat [אני] > [הגדרות] > [כללי] > [שטח אחסון], אפשר לנקות לפי צ'אט וסוג קובץ; בטוח יותר מכלי צד שלישי ל"ניקוי WeChat".
  \item עקרון: מה שיש במערכת/באפליקציה — לא צריך צד שלישי; מה שאפשר לסמן לפי נתיב/סוג — אל תבחרו "ניקוי כולל בלחיצה אחת". מטרת הניקוי היא לשחרר מקום, לא למחוק יחד נתונים חשובים.
\end{itemize}

\booksubsection{תוכנות שימושיות לחסימה ואבטחה}

\begin{itemize}
  \item המרכז הלאומי נגד הונאות קהל יעד: כולם, במיוחד קשישים, עובדי כספים, מי שקונה הרבה ברשת/עונה לשיחות לא מזוהות. נקודותחוזק: רשמי, מאפשר התראות שיחה, אזהרות סיכון SMS, זיהוי כתובות חשודות; מקושר למאגר המאבק בהונאות של המשטרה — החוליה הרשמית החזקה ביותר ב"שלוש שכבות ההגנה".
  \item מרכז אבטחה / מנהל טלפון של המותג קהל יעד: משתמשי המותג המתאים. נקודותחוזק: Huawei Phone Manager, Xiaomi Security Center, OPPO Phone Manager, vivo Phone Manager ועוד — משלבים חסימת שיחות, סריקת וירוסים, ניהול הרשאות והגנת תשלומים, לרוב חינם ומשתפים פעולה טוב יותר עם ליבת המערכת. לא מומלץ להתקין ערכת אבטחה מסחרית נוספת — זה יוצר התנגשויות הרשאות.
  \item שירות חסימת שיחות של הספק קהל יעד: משתמשים שמותקפים בשיחות תדירות גבוהה וקולות AI. נקודותחוזק: ב-China Mobile שולחים KTFSR ל-10086; ב-Telecom יעקבו אחרי "Tianyi Anti-Harassment"; ב-Unicom מגדירים ב-"Unicom Mobile Manager". החסימה מתרחשת בצד הספק — הטלפון לא מצלצל ועדיין נחסמות שיחות זבל.
  \item AdGuard קהל יעד: כל מי שסובל מחלונות קופצים, פרסומות וידאו ופרסומות בתוך אפליקציות. נקודותחוזק: חסימת פרסומות ברמת המערכת (לא רק בדפדפן), סינון באנרים ומעקב באפליקציות, חיסכון בתעבורה והאצת טעינה; הגרסה בתשלום תומכת בכל המכשירים.
  \item היזהרו מהתקנת חבילות "ניקוי/חסימה": חלק מכלי הצד השלישי יושבים ברקע ומבקשים שוב ושוב גישה לאנשי קשר ולהודעות — יכולת החסימה לא תמיד טובה יותר מהמערכת, אבל מוסיפה צריכת סוללה וסיכוני פרטיות. קודם פתחו את הפונקציות המובנות, ואז שקלו אפליקציית מאבק רשמית.
\end{itemize}

\textbf{עקרונות בחירת תוכנות שימושיות}\par

\begin{itemize}
  \item מערכת קודם: חסימה, ניקוי, ניהול הרשאות — קודם פתחו את המתגים המובנים בטלפון.
  \item השלמה רשמית: המרכז הלאומי נגד הונאות ושירותי הספק הם ערוצים רשמיים — אפשר להוסיף.
  \item מתינות בצד שלישי: מספיק להתקין אחד מאותו סוג — להימנע מהתנגשויות הרשאות ומריבוי תהליכים ברקע.
  \item פחות זה יותר: ככל שמותקנות יותר אפליקציות, כך גדלים ההתראות וההרשאות, והסיכוי להפרעה "מהכלי עצמו" גבוה יותר.
\end{itemize}

\booksubsection{חסימת שיחות זבל מובנית בטלפון: איך מפעילים לפי מותג}

להלן כניסות לפי מותג. השמות עשויים להשתנות מעט בין גרסאות — אפשר לחפש בהגדרות “שיחות זבל”, “חסימה”, “שיחות לא מזוהות”. משלים לפרק 8 “טיפול בשיחות זבל”: כאן הדגש על מיקום המתגים, שם על סימון מספרים וערעורים.

\begin{center}
\begin{tabularx}{\linewidth}{|X|X|X|}
\hline
\textbf{פלטפורמה} & \textbf{נתיב כניסה} & \textbf{מומלץ להפעיל} \tabularnewline
\hline
iPhone / iOS & {}[הגדרות] > [טלפון] / [הודעות] & השתקת שיחות לא מזוהות; אנשי קשר חסומים; "סינון שולחים לא מזוהים" בהודעות \tabularnewline
Xiaomi / Redmi & אפליקציית טלפון > הגדרות > חסימת שיחות; או מרכז אבטחה & חסימה חכמה, ניתוק שיחה אחת ממספר לא מזוהה, סינון SMS לפי מילות מפתח \tabularnewline
Huawei / Honor & טלפון > חסימת שיחות; מנהל טלפון > חסימת שיחות & זיהוי מתקשר, רשימה שחורה, חסימת SMS חכמה \tabularnewline
OPPO / realme / OnePlus & אפליקציית טלפון > חסימת שיחות; או מכשיר זה > חסימת שיחות & חסימת שיחות, זיהוי שיחות הונאה, כללי חסימת SMS \tabularnewline
vivo / iQOO & טלפון > חסימת שיחות זבל; iManager & חסימת שיחות, אזהרת שיחות הונאה, סינון הודעות \tabularnewline
Samsung & טלפון > הגדרות > חסימת מספרים / זיהוי מתקשר ושיחות זבל & חסימת מספרים לא מזוהים, זיהוי שיחות זבל (תמיכה משתנה לפי אזור) \tabularnewline
\hline
\end{tabularx}
\end{center}

\begin{tcbinfo}{המלצה לשילוב שלוש שכבות הגנה}
ליום-יום מומלץ “חסימת מערכת + חסימת ספק + המרכז הלאומי נגד הונאות”: המערכת מטפלת בסימון וברשימה שחורה, הספק חוסם שיחות תדירות ברשת, והמרכז מוסיף התראות רשמיות. לא צריך להפעיל הכול ברמה המחמירה ביותר — כדי לא לחסום בטעות שליחויות, משלוחי אוכל ורופאים — אפשר לאשר לאנשי קשר ולאנשי קשר אחרונים.

\end{tcbinfo}

\begin{tcbpractice}{תרגול מעשי: הפעילו חסימת שיחות זבל עוד היום}
\begin{enumerate}
  \item פתחו הגדרות של אפליקציית הטלפון, מצאו "חסימת שיחות / השתקת שיחות לא מזוהות", והפעילו חסימה חכמה.
  \item בדקו גם בטלפון של ההורים: אותו נתיב + הפעלת חסימת הספק (ב-China Mobile אפשר לשלוח KTFSR ל-10086).
  \item התקינו מחנות האפליקציות "המרכז הלאומי נגד הונאות", הירשמו והפעילו התראות שיחה.
  \item נסו להתקשר למספר בדיקה שסומן (או צפו ברשימת החסומים) כדי לוודא שהפונקציה עובדת.
\end{enumerate}

\end{tcbpractice}

\booksection{התראות מזג אוויר, אסונות טבע ורעידות אדמה}

חסימת שיחות זבל עוצרת “הונאות טלפון”, והתראות עוצרות “אסונות טבע ומזג אוויר קיצוני”. אנשים רבים מתקינים אפליקציית מזג אוויר רק לראות טמפרטורה, אבל לא פותחים התראות ולא מפעילים התראת רעידות אדמה — בשניות עד עשרות שניות לפני גשם כבד, טייפון או רעידת אדמה, הטלפון יכול להשיג לך זמן יקר להצלה.

\booksubsection{התראות מזג אוויר: איך מקבלים על גשם כבד, טייפון וחום}

\begin{itemize}
  \item ערוצים רשמיים קודם: ההתראות של לשכת המטאורולוגיה הסינית / התחזית המרכזית / לשכות המחוזיות — הכי סמכותיות. אפשר לעקוב אחרי חשבון WeChat "מזג האוויר של סין", Weibo, או להוריד אפליקציות רשמיות כמו "מזג האוויר של סין".
  \item אפליקציית מזג אוויר של המערכת: Huawei Weather, Xiaomi Weather, OPPO Weather, vivo Weather, "מזג אוויר" של Apple ועוד — רובן דוחפות התראות של לשכת המטאורולוגיה המקומית; היכנסו להגדרות האפליקציה ואשרו שהתראות "מזג אוויר / מזג אוויר אסון" מופעלות.
  \item Caiyun Weather ותחזיות גשם ברמת דקות: מתאים לפני נסיעה לעבודה או איסוף ילדים — "לקחת מטריה בחצי שעה הקרובה?"; יכולת ההתראה נשארת לפי לשכת המטאורולוגיה הרשמית.
  \item משמעות צבעי התראה (כללי): כחול: לשים לב, בדרך כלל לצאת כרגיל ולעקוב אחרי עדכונים.
  \item צהוב: לחזק זהירות, לצמצם פעילות חוץ מיותרת.
  \item כתום: חמור, להישאר בפנים ככל האפשר, להתרחק מאזורי סיכון.
  \item אדום: חמור במיוחד, לפעול לפי הנחיות ביטול לימודים/עבודה ופינוי — לא להיכנס לכבישים מוצפים, שטפונות או רוחות חזקות.
\end{itemize}

סוגי התראות נפוצים: גשם כבד, שלג כבד, טייפון, חום גבוה, גל קור, רוחות חזקות, ברד, קרח בכביש, ערפל כבד, בצורת, סיכון אסון גיאולוגי.

\booksubsection{התראות אסונות טבע: לא רק “גשם”}

\begin{itemize}
  \item מי מפרסם: לשכת המטאורולוגיה אחראית למזג אוויר; אסונות גיאולוגיים ושטפונות מפורסמים לרוב בשיתוף גופי משאבי טבע, מים וחירום; גם ממשלות מקומיות מפרסמות דרך רדיו חירום, SMS, 12379 ועוד.
  \item איך לקבל מכמה ערוצים: התראות מערכת + התראות באפליקציית מזג אוויר.
  \item עקבו ב-WeChat אחרי חשבונות רשמיים של "מנהל החירום" ו"לשכת המטאורולוגיה" המקומית.
  \item ערוצים של 12379 להתראות לאומיות (SMS/שירות, לפי פתיחת אזור).
  \item טלוויזיה, רדיו חירום, רמקולים בכפר/בקהילה — בהרים ובעונת הגשמים אל תסתמכו רק על הטלפון.
\end{itemize}

תזכורות חוץ ונסיעות: ביום טייפון לא להיצמד לשלטי חוצות ולא ללכת תחת עצים; בגשם כבד לא להיכנס לחניונים תת-קרקעיים, מנהרות וגשרים; בטיול בהרים עם התראת "סיכון אסון גיאולוגי" — לצאת מיד ממדרונות ועמקים.

שילוב עם מדריך עזרה ראשונה: אחרי אסון, ניתוק קשר, סוללה ריקה או צורך במיקום לעזרה — ראו פרק 32 רעידות אדמה/אסונות, ופרק 7 תקשורת לווינית ורשת ללא קליטה.

\booksubsection{התראת רעידות אדמה: לא אותו דבר כמו “חיזוי רעידת אדמה”}

התראת רעידת אדמה (EEW) שולחת אזהרה לאזור שעלול להיפגע לפני שגלי הרעידה (בעיקר גל ה-S ההרסני) מגיעים — בזכות כך שגלים אלקטרומגנטיים מהירים יותר מגלי רעידה, ותחנות מזהות את גל ה-P מוקדם יותר. היא לא מנבאת אם תהיה רעידה, אבל משיגה חלון של שניות עד עשרות שניות “להתגונן במקום, לכבות אש, לצאת מהמעלית”. ככל שקרובים יותר למוקד, זמן ההתראה קצר יותר — ולפעמים לא מספיק.

\begin{itemize}
  \item פתרונות רשמיים ומקצועיים: בסין נפוצים שירותי התראה מרשת ההתראות הסינית, וממכון צ'נגדו להפחתת אסונות; חלק מהיצרנים מתקינים מראש או מאפשרים קבלת התראות דרך רשת.
  \item אפליקציות צד שלישי: חפשו בחנות האפליקציות "התראת רעידות אדמה" וכו' (לפני ההתקנה בדקו את המפתח — הימנעו מאפליקציות מזויפות שמבקשות המון הרשאות לא רלוונטיות).
  \item מה לעשות אחרי קבלת אזהרה: בבית להתגונן במקום קרוב (מתחת לשולחן, פינת קיר נושא, ליד רהיט יציב), להגן על הראש והצוואר; להתרחק מחלונות, קירות חיצוניים ורהיטים גבוהים; לא לרוץ במדרגות/לקפוץ; במעלית לרדת מיד בקומה הקרובה; בחוץ לשטח פתוח, להתרחק מקווי חשמל, שלטי חוצות ומבנים מסוכנים. אחרי הרעידה לבדוק גז וחשמל, לשים לב לרעידות משנה.
  \item לא לשתוק: התראת רעידת אדמה צריכה לצלצל "דרך השתקה/לא להפריע"; לפני שינה שימו את הטלפון בהישג יד, ואל תכוונו צליל ל-0.
\end{itemize}

\booksubsection{התראת רעידות אדמה מובנית: איך מפעילים לפי מותג}

להלן כניסות עשויות להשתנות מעט לפי גרסה — אפשר לחפש בהגדרות “רעידת אדמה”, “התראה”, “התראת ביטחון”. ב-iOS התמיכה בהתראות רעידה בסין מוגבלת בדרך כלל — צריך להתקין אפליקציית התראה אמינה, ולאשר מיקום והתראות.

\begin{center}
\begin{tabularx}{\linewidth}{|X|X|X|}
\hline
\textbf{פלטפורמה} & \textbf{כניסה נפוצה} & \textbf{הערה} \tabularnewline
\hline
Xiaomi / Redmi & מנהל טלפון / מרכז אבטחה > התראת רעידות אדמה; או חיפוש "רעידת אדמה" בהגדרות & בחלק מהגרסאות שירות התראה מובנה, דורש חיבור רשת והפעלת מיקום \tabularnewline
Huawei / Honor & הגדרות > אבטחה > התראת חירום / התראת רעידות אדמה & השם עשוי להיות "התראת חירום", כולל רעידות אדמה \tabularnewline
OPPO / realme / OnePlus & הגדרות > אבטחה (SOS) > התראת רעידות אדמה & תלוי בדגם ובכיסוי השירות באזור \tabularnewline
vivo / iQOO & אפליקציית מזג אוויר או הגדרות > אבטחה/עוד > התראת רעידות אדמה & לפי תפריט גרסת המערכת בפועל \tabularnewline
iPhone / iOS & בדרך כלל אין מתג מובנה מלא; התקינו אפליקציית התראה אמינה ואשרו התראות & האפליקציה צריכה לפעול ברקע והרשאות מיקום; אל תחסמו את ההתראות שלה ב"לא להפריע" \tabularnewline
Samsung ועוד & הגדרות > אבטחה וחירום / התראות חירום אלחוטיות & כיסוי "התראות חירום" משתנה לפי אזור וספק \tabularnewline
\hline
\end{tabularx}
\end{center}

\begin{tcbwarn}{טעויות נפוצות על התראת רעידות אדמה}
\begin{itemize}
  \item התראה $\neq$ תחזית: היא לא אומרת "האם תהיה רעידה מחר" — רק כשכבר קרתה והגלים עדיין לא הגיעו.
  \item לא בכל מקום יש הרבה שניות מראש: סמוך למוקד יכולות להיות רק שניות בודדות, ההתראה מועילה בעיקר לאזורים סביב.
  \item כיבוי התראות = כאילו לא מותקן: גם התראה טובה לא מגיעה אם מושתקים, הורגים תהליכים ברקע, או לא מאשרים מיקום.
  \item אפליקציות מזויפות: היזהרו מתוכנות "חיזוי רעידות" שמבקשות גישה לאנשי קשר והודעות — העדיפו מובנה במערכת או פתרון רשמי.
\end{itemize}

\end{tcbwarn}

\begin{tcbinfo}{איך לשלב יכולות התראה}
המלצה ליומיום: “התראות מזג אוויר של המערכת + ערוצים רשמיים/חירום + התראת רעידות אדמה (מערכת או אפליקציה רשמית)”. כמו בחסימה — לא צריך הרבה תוכנות מאותו סוג; המפתח הוא אישור התראות, מיקום, ויכולת לחדור “לא להפריע”. אחרי שחסימת השיחות מופעלת בטלפון של ההורים, בדקו גם התראות מזג אוויר ורעידות אדמה.

\end{tcbinfo}

\begin{tcbpractice}{תרגול מעשי: בדיקת הגדרות התראה ב-10 דקות}
\begin{enumerate}
  \item פתחו אפליקציית מזג אוויר של המערכת, היכנסו להגדרות, הפעילו התראות "מזג אוויר / מזג אוויר אסון".
  \item בהגדרות המערכת חפשו "רעידת אדמה" או "התראה" והפעילו התראת רעידות/חירום (Android); ב-iPhone התקינו אפליקציית התראה רשמית ואשרו התראות.
  \item עקבו אחרי החשבון הרשמי של לשכת המטאורולוגיה/החירום המקומית ואשרו דחיפה.
  \item ודאו שצליל ההתראה לא מושתק לחלוטין ב"לא להפריע"; בדקו אותו דבר בטלפון של ההורים.
\end{enumerate}

\end{tcbpractice}

\booksection{שיתוף בין מכשירים: שליטה מרחוק, העברת קבצים ואבטחת חשבונות}

טלפון שעובד חלק לרוב “עובד בצוות” עם מחשב וטאבלט. בסעיף זה נוסיף כלים שימושיים בתדירות גבוהה: שליטה מהטלפון במחשב, עבודה עם AI ושיתוף מרחוק, טפטים חיים, העברת קבצים ברשת מקומית, לוח-גזירים בין מכשירים, אסימוני 2FA, וחיבור סריקת נייר ללוחות זמנים ומשימות. לשילוב מכשירים שיטתי יותר ראו פרק 17.

\booksubsection{שליטה מהטלפון במחשב: Sunlogin / ToDesk}

\begin{itemize}
  \item Sunlogin Remote Control קהל יעד: מי שצריך לתפעל Windows או Mac בבית/בעבודה כשלא נמצא; גם נפוץ לתיקון מרחוק של טלפונים להורים (ראו פרק 27). נקודותחוזק: אקוסיסטם סיני בוגר, אפליקציית טלפון שולטת בשולחן העבודה, מעבירה קבצים, מדליקה/מכבה (חלק דורש חומרה); הגרסה החינמית מספיקה לרוב השימוש האישי.
  \item ToDesk קהל יעד: מי שמחפש חיבור יציב וזמן השהיה נמוך לעבודה/תחזוקה מרחוק. נקודותחוזק: חיבור עם קוד מכשיר + סיסמה, שליטה מהטלפון במחשב, מהמחשב בטלפון (לפי גרסה), העברת קבצים והחלפת מסכים מרובים — חוויית צמתים בסין לרוב טובה.
  \item שימוש בסיסי: התקינו לקוח במחשב והתחברו/הישארו מחוברים התקינו אפליקציה בטלפון התחברו עם מזהה מכשיר + סיסמת חיבור (או אותו חשבון) תפעלו את שולחן העבודה עם משטח מגע/מחוות.
  \item קווים אדומים לאבטחה (חובה לזכור): זר שמציג עצמו כ"שירות/רשויות/ספק" ומבקש שתתקינו Sunlogin/ToDesk ושליטה מרחוק — דחו תמיד — זו הונאה טיפוסית (ראו פרק 27).
  \item בשימוש עצמי: הפעילו אימות דו-שלבי, שנו סיסמת ברירת מחדל, התנתקו או כבו גישה ללא השגחה כשלא בשימוש; במחשב עבודה הקפידו על כללי האבטחה של הארגון.
  \item במקום ציבורי בעבודה מרחוק — שימו לב אם מישהו מציץ במסך או בסיסמה.
\end{itemize}

\booksubsection{תוכנות משרד AI: Trae / WorkBuddy / Dumate}

\begin{itemize}
  \item Trae קהל יעד: משתמשי תכנות ושיתוף מסמכים, ידידותי לסביבה סינית. נקודותחוזק: IDE/עוזר תכנות AI מבית ByteDance, מצב SOLO יוצר פרויקט במשפט אחד; יעיל לכתיבת קוד ועריכת מסמכים. אם רץ על מחשב, אפשר דרך שליטה מרחוק או יכולות רשמיות לעקוב במכשיר. ראו פרק 41 נוף כלי תכנות AI.
  \item WorkBuddy קהל יעד: עובדי משרד שרוצים AI לסדר סיכומי פגישות, מיילים, דוחות ומסמכים. נקודותחוזק: עוזר AI לתרחישי משרד, מתאים ל"טיוטה — ליטוש — סיכום"; חלק מזרימות העבודה אפשר להתחיל בטלפון, והעבודה הכבדה במחשב.
  \item Dumate — למי מתאים: לצוותים וליחידים שזקוקים לשיתוף פעולה רב-משתתפים ולניהול משולב של מסמכים וידע. נקודות החוזקה העיקריות: כלי עבודה משרדי בגישת AI, עם דגש על ריכוז חומרים ועריכה משותפת; בנייד אפשר לצפות ולבצע עריכה קלה, ועימוד מורכב עדיף לעשות על המסך הגדול של המחשב.
  \item איך להבין "שליטה מהטלפון במחשב": הרבה יכולות משרד AI רצות על מחשב/דפדפן; בטלפון דרך Sunlogin/ToDesk מתחברים למחשב, או משתמשים באפליקציה/אתר של המוצר, וכך "משדרים" משימות AI בנסיעה. עיקרון: חישוב כבד ומסמכים ארוכים במחשב; צפייה בסטטוס ועריכות קצרות בטלפון.
  \item תזכורת פרטיות: לפני העלאת חוזים, תעודות זהות ונתוני לקוחות — קראו את מדיניות הפרטיות ודרישות התאימות; לחומרי עבודה העדיפו את הכלי שהארגון מגדיר.
\end{itemize}

\booksubsection{שולחן עבודה מלאכותי: טפטים חיים של miHoYo (טלפון + PC)}

\begin{itemize}
  \item Artificial Desktop קהל יעד: אוהבי דמויות miHoYo שרוצים טפטים חיים חינם. נקודותחוזק: אפליקציית טפטים חיים חינם מבית miHoYo, תומכת בטלפון ו-PC; דמויות Live2D/הופעות חיות יכולות לשמש כשולחן עבודה, בחלק מהמכשירים/המערכות אפשר להגדיר כמסך נעילה או טפט.
  \item סוללה וביצועים: טפט חי צורך יותר סוללה ו-GPU מתמונה סטטית; בסוללה חלשה או מכשיר ישן אפשר להשתמש רק בטעינה, או להוריד קצב פריימים (אם האפליקציה מציעה).
  \item מקורות התקנה: העדיפו את האתר הרשמי או חנויות אפליקציות מוכרות, והימנעו מ"גרסאות פרוצות" של צד שלישי המגיעות עם סוסים טרויאניים; בגרסת המחשב שימו לב לנתיב ההתקנה ולשאלה אם בכלל צריך הפעלה אוטומטית עם אתחול המחשב.
\end{itemize}

\booksubsection{העברת קבצים ברשת מקומית: LocalSend / Feige / AirDrop}

\begin{itemize}
  \item LocalSend קהל יעד: מי שמערב Android/Windows/iOS/Mac/Linux ורוצה העברה בין פלטפורמות בלי ענן ציבורי. נקודותחוזק: קוד פתוח, חיבור ישיר ברשת מקומית, העברת תמונות ווידאו ומסמכים באותו Wi-Fi, בלי כבל ולא דרך שרתי יצרן — המהירות תלויה בנתב.
  \item Feige — למי מתאים: לרשתות של מחשבי Windows במשרד ולסביבות עבודה על רשת מקומית ותיקה. נקודות החוזקה העיקריות: כלי ותיק להעברת קבצים וצ'אט ברשת מקומית, מתאים להעברת קבצים גדולים בתוך הרשת הפנימית; הורידו רק מערוץ רשמי, וברשת הארגונית פעלו לפי מדיניות ה-IT.
  \item AirDrop קהל יעד: משתמשי iPhone, iPad, Mac. נקודותחוזק: שיתוף בלחיצה אחת ב-Apple, גילוי Bluetooth + Wi-Fi ישיר, בלי תעבורה; אפשר להגדיר "אנשי קשר בלבד" כדי להפחית הטרדות.
  \item תוספות: Xiaomi Share, Huawei Share, vivo Share ועוד פתרונות מותג — מהירים בין מכשירים מאותו מותג; "מסייע העברת קבצים" ב-WeChat נוח אבל מוגבל בגודל קובץ ומכשירים — לקבצים גדולים העדיפו LocalSend/AirDrop.
  \item טיפים: שני הצדדים לאותו Wi-Fi; רשת Wi-Fi ציבורית להעברת קבצים רגישים מסוכנת — העדיפו נקודה אישית; אחרי ההעברה ודאו בצד המקבל שהקובץ שלם.
\end{itemize}

\booksubsection{סנכרון לוח-גזירים בין מכשירים: KDE Connect / שיתוף טלפון-מחשב}

\begin{itemize}
  \item "לוח גזירים אוניברסלי" של Apple קהל יעד: משתמשי iPhone + Mac. נקודותחוזק: אחרי התחברות לאותו Apple ID והפעלת Handoff, מעתיקים בטלפון מדביקים ב-Mac (ולהפך), טקסט וחלק מהקבצים.
  \item KDE Connect קהל יעד: Android + Windows/Linux/Mac, מי שרוצה סנכרון לוח-גזירים בקוד פתוח, שיקוף התראות וקלט מרחוק. נקודותחוזק: אחרי צימוד ברשת מקומית — סנכרון לוח-גזירים, דחיפת התראות מהטלפון למחשב, שימוש בטלפון כשלט/משטח מגע; הנתונים לא עוברים בחשבון צד שלישי.
  \item "שיתוף טלפון-מחשב" של מותגים: Huawei Multi-Screen Collaboration, Xiaomi MiMo, vivo Share, OPPO Connect, Honor Connect ועוד — לרוב עם לוח-גזירים וגרירת קבצים בין מסכים, חווייה מיטבית באותו מותג.
  \item ערוצים רשמיים של Windows + Android: "Phone Link / Link to Windows" (התפתחות של Your Phone) — במחשב רואים התראות ותמונות מהטלפון, ותומך בהעתקה-הדבקה בתרחישים נתמכים; ב-Samsung האינטגרציה עמוקה יותר.
  \item תזכורת אבטחה: לוח-גזירים עלול להכיל סיסמאות וקודי אימות; במחשב משותף היזהרו מסנכרון גלובלי, והימנעו מהדבקה חוזרת בשדה סיסמה; חלק ממקלדות גם קוראות לוח-גזירים — שימו לב להרשאות.
\end{itemize}

\booksubsection{אסימוני 2FA: Aegis / Microsoft Authenticator}

\begin{itemize}
  \item למה להשתמש במאשר ולא ב-SMS: קוד SMS עלול להיחטף בהחלפת SIM; אסימון TOTP נוצר מקומית לפי אלגוריתם זמן, בלי תלות בסלולר — עמיד יותר להונאה וליירוט (עדיין צריך להיזהר מסוסים טרויאניים וצילומי מסך).
  \item Aegis Authenticator קהל יעד: משתמשי Android שמעדיפים קוד פתוח ושליטה מקומית. נקודותחוזק: חינם וקוד פתוח, גיבוי מוצפן וייבוא/ייצוא, ממשק נקי, נעילת רשימת אסימונים עם טביעת אצבע/פנים.
  \item Microsoft Authenticator קהל יעד: משתמשי חשבון Microsoft, Office 365/חשבונות ארגוניים; עובד בין iOS/Android. נקודותחוזק: אישור כניסה בדחיפה, ניטור סיסמאות וגיבוי חשבון (לפי גרסת מוצר), אינטגרציה הדוקה עם Windows/זהות ארגונית.
  \item Google Authenticator / 2FAS ועוד: נפוצים גם בתוך האקוסיסטם; העיקר שיש גיבוי מוצפן — במעבר מכשיר לא מאבדים אסימונים.
  \item חובה במעבר מכשיר: העברה מהמכשיר הישן או שימוש ב"קודי גיבוי/מפתח שחזור" מהשירות; הדפיסו קודי גיבוי על נייר ושמרו בשני מקומות — לא רק באותו טלפון.
  \item עם מנהל סיסמאות: חלק מכספות הסיסמאות גם מאחסנות TOTP — נוח אבל מרוכז סיכון; לחשבונות בעלי ערך מומלץ "מנהל סיסמאות + מאשר עצמאי". ראו פרק 14.
\end{itemize}

\booksubsection{שילוב מומלץ: איך מחברים יום שלם}

\textbf{דוגמת זרימת עבודה שימושית}\par

\begin{itemize}
  \item בוקר: TickTick ללוח הזמנים של היום מבט על מזג אוויר/התראות פותחים Trae/מסמכים במחשב, הטלפון בצד.
  \item יום: חומרי נייר סורקים ל-PDF, שולחים למחשב דרך LocalSend/AirDrop; מסמכים ארוכים מנסחים במחשב עם כלי משרד AI.
  \item בחוץ: מתחברים חזרה למחשב דרך ToDesk/Sunlogin וממשיכים במשימה; כניסה לחשבונות חשובים עם אסימון דינמי של Aegis/Microsoft Authenticator.
  \item ערב: מסמנים ביצועים ב-TickTick; אם צריך, מחליפים טפט ב-Artificial Desktop להירגעות.
\end{itemize}

\begin{tcbwarn}{כללי אצבע לשליטה מרחוק ואבטחת חשבונות}
\begin{itemize}
  \item מישהו מבקש שתתקינו תוכנת שליטה מרחוק ניתוק שיחה, התקשרו בעצמכם לשירות הרשמי.
  \item קודי אימות ואסימונים דינמיים לא נמסרים לאף אחד, כולל "שליח" ו"שירות לקוחות".
  \item לפני שליטה מרחוק במחשב עבודה או התקנת תוכנת AI — שאלו את ה-IT אם מותר.
  \item ברשת ציבורית — זהירות יתרה בהעברת קבצים וכניסה מרחוק.
\end{itemize}

\end{tcbwarn}

\begin{tcbpractice}{תרגול מעשי: הפעילו את שרשרת "בין מכשירים" פעם אחת}
\begin{enumerate}
  \item חברו טלפון ומחשב לאותו Wi-Fi, העבירו תמונה עם LocalSend (או AirDrop).
  \item בנו משימה ב-TickTick "היום 20:00: לארגן סריקות" וודאו שרואים גם בטלפון וגם במחשב.
  \item הפעילו 2FA למייל או ל-WeChat (אם נתמך), קשרו עם Aegis או Microsoft Authenticator, ושמרו קודי גיבוי.
  \item התקינו Sunlogin או ToDesk במחשב שלכם, התחברו פעם אחת מהטלפון, ואז התנתקו או חיזקו סיסמה — להכיר את התהליך בלי להשאיר דלת אחורית.
\end{enumerate}

\end{tcbpractice}
